%% kapa: exam=1 round=34 year=2023 period=2 form=A subject=accounting subject_ko=회계학 seq=40 number=80
%% source: Q-Net 감정평가사 기출문제 (cst00309) artlSeq=5211789 file=raw/1008/5211789/2023년도 제34회 감정평가사 1차시험 2교시.pdf
%% conversion: kordoc 4.24.51 -> tools/gen1.py (auto); review=proofread-sonnet; flags=none
\begin{question}{80}
(주)감평은 가중평균법에 의한 종합원가계산제도를 채택하고 있으며, 단일공정을 통해 제품을 생산한다. 모든 원가는 공정 전반에 걸쳐 균등하게 발생한다. (주)감평의 당기 생산 관련 자료는 다음과 같다.
\par{\small\centering\begin{tabular}{|c|c|c|c|}\hline
구 분 & 물량(완성도) & 직접재료원가 & 전환원가 \\ \hline
기초재공품 & 100단위 (?) & ￦4,300 & ￦8,200 \\ \hline
당기착수 & 900 & 20,000 & 39,500 \\ \hline
기말재공품 & 200 (?) & ? & ? \\ \hline
\end{tabular}\par}
(주)감평의 당기 완성품환산량 단위당 원가가 ￦80이고 당기 완성품환산량이 선입선출법에 의한 완성품환산량보다 50단위가 더 많을 경우, 선입선출법에 의한 기말재공품 원가는? (단, 공손 및 감손은 발생하지 않는다.)
\begin{choices}
  \item ￦3,500
  \item ￦4,500
  \item ￦5,500
  \item ￦6,500
  \item ￦7,000
\end{choices}
\end{question}
\answer{5}
